% SAMPLE-ONLY: \SampleGuideNote and \SampleGuideSource are not template commands.
\section{考察}
\label{sec:discussion}

\subsection{データ処理}
\label{subsec:data-processing}

CSV ファイルから周波数、利得、位相差を読み込み、理論値との差を計算した。
\Cref{lst:analysis}に、読み込みと基本計算の処理を示す。

% [DOC:CODE-BLOCK-01] Python listing with comments, strings, numbers, and wrapping.
% guide_listing_start
\begin{lstlisting}[
  language=Python,
  caption={周波数応答データの読み込みと基本計算},
  label={lst:analysis}
]
import csv
import math

INPUT_FILE = "data/frequency-response.csv"

with open(INPUT_FILE, encoding="utf-8") as stream:
    rows = list(csv.DictReader(stream))

# Convert selected CSV fields from text to floating-point values.
measurements = [
    {
        "frequency": float(row["frequency_hz"]),
        "gain": float(row["gain_measured"]),
        "phase": float(row["phase_measured"]),
    }
    for row in rows
]

gain_cutoff = min(
    measurements,
    key=lambda row: abs(row["gain"] + 3.0),
)

message = (
    f"nearest cutoff candidate: {gain_cutoff['frequency']:.0f} Hz, "
    f"gain={gain_cutoff['gain']:.2f} dB, phase={gain_cutoff['phase']:.1f} deg"
)
print(message)
\end{lstlisting}
% guide_listing_end

\SampleGuideSource
  {コードを本文に載せるための LaTeX ソース}
  {sections/05-discussion.tex}
  {guide_listing_start}
  {guide_listing_end}

\SampleGuideNote[red!65!black]{［CODE-BLOCK-01］}{%
\texttt{lstlisting} の中にプログラムを書くと、コード用の書体と枠で表示できます。
\texttt{caption} には見出し、\texttt{label} には参照用の名前を書きます。}

% [DOC:CODE-HIGHLIGHT-01] Syntax highlighting by language.
\SampleGuideNote[red!65!black]{［CODE-HIGHLIGHT-01］}{%
\texttt{language=Python} により、予約語、コメント、文字列が種類ごとに色分けされます。}

% [DOC:CODE-WRAP-01] Automatic wrapping of long source lines.
\SampleGuideNote[red!65!black]{長い行の補足}{%
枠幅を超えるコードは、全体設定の \texttt{breaklines=true} により次の行へ折り返されます。}

% [DOC:CODE-REF-01] Cross-reference to a listing label.
\SampleGuideNote[red!65!black]{コード参照の補足}{%
\texttt{caption} と \texttt{label} を付けたコードは、図や表と同様に \texttt{\string\cref} から参照できます。}

\subsection{利得から求めた遮断周波数}
\label{subsec:gain-cutoff}

$-3\,\mathrm{dB}$ をはさむ 1.50 kHz と 1.60 kHz の二点を用い、
利得を周波数に対して線形補間した。補間処理を\cref{lst:interpolation}に示す。

% [DOC:CODE-BLOCK-02] A second Python listing.
\begin{lstlisting}[
  language=Python,
  caption={二点間の線形補間},
  label={lst:interpolation}
]
def interpolate_x(x1, y1, x2, y2, target_y):
    ratio = (target_y - y1) / (y2 - y1)
    return x1 + ratio * (x2 - x1)

cutoff_gain_hz = interpolate_x(
    1500.0, -2.78,
    1600.0, -3.10,
    -3.00,
)

print(f"gain estimate = {cutoff_gain_hz:.1f} Hz")
\end{lstlisting}

得られた遮断周波数は $1.57\,\mathrm{kHz}$ であった。
測定点を一つ選ぶだけの場合は 1.60 kHz となるが、二点補間により理論値へ近い値が得られた。

\subsection{位相差から求めた遮断周波数}
\label{subsec:phase-cutoff}

位相差が $-45\,\mathrm{deg}$ となる点についても、1.50 kHz の
$-44\,\mathrm{deg}$ と 1.60 kHz の $-46\,\mathrm{deg}$ の間を補間した。
この方法による遮断周波数は $1.55\,\mathrm{kHz}$ となった。

利得による推定値と位相差による推定値の差は 20 Hz であった。
この差は、振幅と時間差を別々に読み取ったこと、および表示分解能が異なることにより生じたと考えられる。

\subsection{曲線全体を用いた推定}
\label{subsec:curve-fit}

全測定点に対して理論式の $f_c$ だけを変化させ、利得残差の二乗和が最小となる値を探索した。
実行時に用いたコマンドを\cref{lst:run-command}に示す。

% [DOC:CODE-LANGUAGE-02] Bash listing.
\begin{lstlisting}[
  language=bash,
  caption={解析プログラムの実行コマンド},
  label={lst:run-command}
]
python3 analyze_response.py \
  --input data/frequency-response.csv \
  --fit gain \
  --output build/fit-summary.txt
\end{lstlisting}

\SampleGuideNote[red!65!black]{［CODE-LANGUAGE-02］}{%
\texttt{language} の値を変更すると、同じ \texttt{lstlisting} で別の言語の構文を表示できます。}

曲線全体から得られた遮断周波数は $1.58\,\mathrm{kHz}$ であった。
三つの推定結果を\cref{tab:cutoff-estimates}にまとめる。

\begin{table}[H]
  \centering
  \caption{遮断周波数の推定結果}
  \label{tab:cutoff-estimates}
  \begin{tabularx}{0.88\linewidth}{
      >{\raggedright\arraybackslash}p{0.26\linewidth}
      >{\centering\arraybackslash}p{0.18\linewidth}
      >{\raggedright\arraybackslash}X}
    \toprule
    推定方法 & 推定値 / kHz & 使用した情報 \\
    \midrule
    利得の二点補間 & 1.57 & $-3\,\mathrm{dB}$ をはさむ 1.50 kHz と 1.60 kHz の測定値 \\
    位相差の二点補間 & 1.55 & $-45\,\mathrm{deg}$ をはさむ二つの位相差 \\
    利得曲線の当てはめ & 1.58 & 100 Hz から 10 kHz までの全利得測定値 \\
    \bottomrule
  \end{tabularx}
\end{table}

三つの値は 1.55--1.58 kHz の範囲にあり、素子の測定値から求めた
理論値 $1.58\,\mathrm{kHz}$ と近かった。

\subsection{理論値との差}
\label{subsec:relative-error}

代表値として曲線当てはめによる $1.58\,\mathrm{kHz}$ を用い、
公称値から求めた $1.59\,\mathrm{kHz}$ との相対差を\cref{eq:relative-error}で計算した。

\begin{equation}
  \varepsilon
  =\frac{\lvert 1.58-1.59\rvert}{1.59}\times100
  \simeq0.63\,\%
  \label{eq:relative-error}
\end{equation}

一方、実測した $R=0.998\,\mathrm{k}\Omega$、$C=0.101\,\mu\mathrm{F}$ を
\cref{eq:cutoff}に代入した値は $1.58\,\mathrm{kHz}$ であり、曲線当てはめの値と一致した。

\subsection{素子公差の影響}
\label{subsec:tolerance-discussion}

\crefrange{eq:cutoff-min}{eq:cutoff-max}で求めた 1.49--1.71 kHz の範囲には、
三つの推定値がすべて含まれる。一次近似では、遮断周波数の相対変化は
抵抗と静電容量の相対変化の和に対応する。

\begin{align}
  \frac{\Delta f_c}{f_c}
    &\simeq-\frac{\Delta R}{R}-\frac{\Delta C}{C}
      \label{eq:tolerance-propagation}\\
  \left\lvert\frac{\Delta f_c}{f_c}\right\rvert
    &\leq
      \left\lvert\frac{\Delta R}{R}\right\rvert
      +\left\lvert\frac{\Delta C}{C}\right\rvert.
      \label{eq:tolerance-bound}
\end{align}

測定した抵抗値のずれは小さく、今回の遮断周波数には主に静電容量のずれが影響したと考えられる。
素子公差を含む測定不確かさの整理には、学位論文で示された方法も参照した\cite{jp-thesis}。

\subsection{測定器による負荷}
\label{subsec:instrument-loading}

理論式では信号発生器を理想電圧源、オシロスコープを無限大の入力抵抗として扱った。
実際には信号発生器に $50\,\Omega$ の出力抵抗があり、プローブにも入力抵抗と入力容量がある。
そのため、回路から見た実効抵抗と実効容量は公称値だけで決まらない。

特に高周波側ではプローブの入力容量がコンデンサと並列に加わり、
理論曲線よりわずかに大きな減衰を生じる可能性がある。
\cref{fig:gain-residual}で 10 kHz の残差が正側へずれたことには、
振幅読取りの分解能や信号発生器の出力変化も影響したと考えられる。

\subsection{走査方向と再現性}
\label{subsec:sweep-direction}

500 Hz から 2.00 kHz までを逆方向にも測定した結果、同じ周波数での利得差は
最大 $0.04\,\mathrm{dB}$ であった。RC 回路は受動回路であり、今回の時間間隔では
明確な履歴依存性は見られなかった。

ただし、周波数を切り替えた直後に値を読み取ると、オシロスコープの自動測定値が
一時的に変動することがあった。測定許可信号の考え方や同期測定については、
講演資料と会議録の整理も参照した\cite{jp-inproceedings,jp-proceedings}。

\subsection{主な誤差要因}
\label{subsec:error-sources}

理論値と測定値の差に含まれる要因を以下にまとめる。

% [DOC:LIST-DESCRIPTION-04] Description list for error sources.
\begin{description}[leftmargin=10em,labelwidth=8.5em,itemsep=0.25em]
  \item[素子の公差] 抵抗値と静電容量が公称値からずれる。
  \item[周波数間隔] 真の $-3\,\mathrm{dB}$ 点を直接測定せず、測定点間を補間している。
  \item[出力抵抗] 信号発生器の $50\,\Omega$ を単純な理論式に含めていない。
  \item[プローブ負荷] プローブの入力抵抗と入力容量が回路へ加わる。
  \item[時間差読取り] カーソル位置の分解能が位相差へ変換される。
  \item[振幅読取り] 高周波側では出力振幅が小さく、相対的な読取り誤差が増える。
\end{description}

実験室内の事前メモに記録された注意点とも比較し、グラウンド接続とプローブ倍率を
毎回確認する必要があると判断した\cite{jp-unpublished}。
英語文献に示された一次フィルタの評価方法とも、測定手順と結果の傾向は一致した\cite{en-article}。

\subsection{測定方法の改善案}
\label{subsec:improvements}

次回の測定では、以下の順に条件を改善する。

% [DOC:LIST-ENUMERATE-02] A second enumerate example.
\begin{enumerate}[label=\arabic*.,leftmargin=3em,itemsep=0.35em]
  \item 1.45--1.70 kHz の測定間隔を 25 Hz 程度まで小さくする。
  \item 各周波数で五回以上測定し、平均値と標準不確かさを求める。
  \item 信号発生器の出力抵抗を含む等価回路で理論値を計算する。
  \item プローブの入力容量を実測または仕様値から補正する。
  \item 自動保存した波形から振幅と位相差を同じ処理で計算する。
\end{enumerate}

\SampleGuideNote[teal!70!black]{番号形式の補足}{%
二つ目の \texttt{enumerate} では \texttt{label=\string\arabic*.} とし、通常の「1.」形式を使用しています。}
